\section{Position du problème}

\subsection{Un contexte d'entreprise marqué par un changement d'activité}

Audit Énergie est une entreprise dont l'activité historique s'est construite autour du
courtage en énergie. Cette activité consiste à accompagner des clients professionnels dans
la compréhension, la comparaison et la renégociation de leurs contrats d'énergie. Dans
un marché devenu progressivement plus complexe, cette activité a conduit l'entreprise à
développer des compétences concrètes dans plusieurs domaines : la prospection, la
relation client à distance, l'analyse d'un besoin commercial, l'argumentation et le suivi
d'un dossier dans la durée.

Ce positionnement initial est important pour comprendre le sujet du mémoire, mais il ne
constitue pas le cœur du travail présenté. Le projet étudié s'inscrit principalement dans
l'évolution d'Audit Énergie vers une activité d'organisme de formation. Cette évolution
s'est faite progressivement. Dans un premier temps, la formation répondait à un besoin
interne : former des profils capables de rejoindre les équipes commerciales ou de relation
client. Cette première étape était cohérente avec l'activité de courtage, car elle permettait
de transmettre des méthodes directement liées au métier de l'entreprise.

Depuis environ un an, cette activité a pris une orientation plus structurée. L'entreprise a
cherché à développer une offre de formation externe, destinée à des apprenants suivant
des parcours professionnalisants à distance. Les formations concernées portent notamment
sur des titres professionnels liés à la relation client à distance ou aux métiers commerciaux,
comme le titre professionnel de conseiller relation client à distance ou celui d'employé
commercial. Dans ce cadre, l'entreprise ne se limite plus à former ses propres équipes :
elle doit produire, organiser, diffuser et suivre des contenus pédagogiques pour des élèves
inscrits dans un parcours de formation.

Ce changement d'échelle modifie profondément les contraintes de l'entreprise. Former
quelques collaborateurs en interne ne demande pas la même organisation que faire
fonctionner une activité de formation externe. Dans le premier cas, la transmission peut
reposer sur l'expérience des salariés, sur des échanges directs ou sur des supports encore
peu formalisés. Dans le second cas, il faut construire des parcours plus reproductibles,
préparer des contenus, assurer une continuité pédagogique, suivre la présence des élèves,
proposer des exercices, répondre aux questions et garantir une qualité suffisante pour que
la formation soit exploitable.

Le sujet du mémoire naît précisément de cette transition. Il ne s'agit pas seulement
d'ajouter une fonctionnalité technique à une plateforme existante, mais de répondre à un
problème d'organisation plus large : comment une petite structure peut-elle développer
une activité de formation à distance sans disposer des mêmes moyens humains, financiers
et organisationnels qu'un organisme de formation plus établi ?

\subsection{La formation à distance comme opportunité et comme contrainte}

La formation à distance représente une opportunité importante pour une petite entreprise.
Elle permet de toucher des apprenants sans dépendre d'une salle physique, de mutualiser
certains contenus, de planifier des sessions plus souples et de réutiliser des supports d'une
promotion à l'autre. Elle semble donc, au premier abord, compatible avec les contraintes
d'une structure légère.

Cependant, cette impression est incomplète. Une formation à distance ne devient rentable
que si les coûts de production et d'animation restent maîtrisés. Or, une formation ne se
résume pas à mettre des documents en ligne. Elle suppose de concevoir un déroulé
pédagogique, de produire les cours, de les rendre accessibles, d'accompagner les
apprenants, de suivre leur participation et de vérifier leur compréhension. Ces tâches sont
répétitives, mais elles restent indispensables.

Dans le cas d'Audit Énergie, le modèle le plus naturel aurait été de faire appel à des
formateurs freelance pour assurer les cours en ligne. Ce modèle présente des avantages
évidents. Un formateur humain peut adapter son discours, répondre aux questions,
reformuler une notion difficile, donner des exemples et maintenir une dynamique de
groupe. Il incarne également une forme de responsabilité pédagogique immédiate : si un
élève ne comprend pas, il peut intervenir directement.

Mais ce modèle pose un problème économique fort. Pour une petite structure, le coût
récurrent d'un formateur peut absorber une part importante de la valeur produite par la
formation. Plus le nombre d'heures est élevé, plus cette contrainte devient structurante.
Dans un marché où les marges doivent rester suffisantes pour financer l'organisation, le
suivi administratif, les outils et le développement commercial, dépendre fortement d'un
formateur externe rend le modèle difficile à industrialiser.

Le problème initial peut donc être formulé simplement : l'entreprise voulait développer
une activité de formation externe, mais le modèle classique reposant sur des formateurs
humains en direct n'était pas compatible avec ses objectifs de coût, de disponibilité et de
passage à l'échelle. Il fallait trouver une manière de produire et diffuser des formations
avec moins d'interventions humaines récurrentes.

Cette contrainte n'est pas propre à Audit Énergie. Elle correspond à une classe de
problèmes fréquente dans les petites entreprises : comment développer une activité de
service exigeante en temps humain, alors que la structure ne dispose pas encore des
ressources nécessaires pour recruter ou mobiliser durablement plusieurs profils
spécialisés ? Dans le cas de la formation, cette question est particulièrement sensible, car
la réduction des coûts ne doit pas conduire à une baisse excessive de la qualité
pédagogique.

\subsection{Une première réponse : dissocier la production du cours et sa diffusion}

La première réponse apportée à cette contrainte a consisté à ne plus raisonner en cours
en direct, mais en cours préparé à l'avance. L'idée était de construire une plateforme
capable d'héberger des contenus audio et de les diffuser automatiquement le jour prévu
pour la formation. Au lieu de rémunérer un formateur pour assurer plusieurs heures de
cours en direct, le travail était séparé en deux étapes : d'abord la préparation du contenu,
puis sa lecture sous forme d'audio.

Ce choix répondait à une logique économique claire. Il devenait possible de faire rédiger
un cours en amont, éventuellement avec l'aide d'outils d'intelligence artificielle, puis de
le faire enregistrer par une personne dont le rôle se limitait principalement à l'interpréter.
La plateforme assurait ensuite la diffusion au moment prévu. Ainsi, une partie du travail
du formateur était transformée en production de contenu réutilisable.

Cette première version de la plateforme reposait donc sur un principe simple : préparer
les cours pendant la semaine, déposer les fichiers audio dans l'outil, puis lancer leur
diffusion le jour de la formation. Pour les élèves, le cours se présentait comme une
session organisée dans le temps. Pour l'entreprise, la production était plus facilement
planifiable qu'un cours en direct assuré par un freelance.

Cette organisation constituait une première étape importante, car elle permettait de
valider plusieurs hypothèses. D'abord, les apprenants pouvaient suivre une formation à
partir de contenus audio diffusés à distance. Ensuite, l'entreprise pouvait internaliser une
partie de la production pédagogique. Enfin, la plateforme devenait un outil central pour
organiser le déroulement de la formation, et non plus seulement un espace de dépôt de
fichiers.

Néanmoins, cette solution ne résolvait qu'une partie du problème. Elle permettait de
réduire la dépendance au formateur en direct, mais elle conservait une forte dépendance
au travail humain. Les cours devaient toujours être rédigés, relus, découpés, enregistrés,
déposés et vérifiés. À chaque nouvelle formation, ou même à chaque mise à jour d'un
contenu, une chaîne de production manuelle devait être relancée.

Le gain économique existait, mais il était limité par le caractère encore artisanal du
processus. Le modèle ne reposait plus sur un formateur freelance animant la formation
en direct, mais il reposait toujours sur des personnes capables de produire et de préparer
les supports à temps. En pratique, l'entreprise avait déplacé le coût humain plutôt qu'elle
ne l'avait supprimé.

\subsection{L'arrivée de l'IA comme changement de perspective}

L'arrivée des outils d'intelligence artificielle générative a modifié la manière de poser le
problème. Dans un premier temps, l'IA pouvait être vue comme une aide à la rédaction.
Un collaborateur pouvait utiliser un modèle de langage pour produire une première
version d'un cours, reformuler un passage ou structurer une séquence pédagogique. Cette
utilisation restait proche d'un usage bureautique assisté : l'humain gardait le contrôle du
processus et utilisait l'IA pour accélérer certaines tâches.

Mais le contexte technologique a rapidement ouvert une hypothèse plus ambitieuse. Si
un modèle de langage peut rédiger un contenu pédagogique, si un outil de synthèse vocale
peut le lire, si un système peut générer des supports visuels et si un assistant
conversationnel peut répondre aux questions des apprenants, alors il devient possible de
concevoir une plateforme beaucoup plus autonome. Dans cette perspective, l'IA ne sert
plus seulement à aider un rédacteur. Elle devient une composante centrale de la chaîne de
production et d'accompagnement pédagogique.

Cette idée doit être formulée sans ambiguïté. L'objectif du projet était bien de remplacer
une partie importante des interventions humaines initialement nécessaires à la formation.
Dans une petite entreprise, ce choix n'est pas seulement technologique : il est économique
et stratégique. L'entreprise cherche à rendre possible une activité de formation qui serait
trop coûteuse si elle reposait entièrement sur des ressources humaines mobilisées en
continu.

Cependant, remplacer une fonction humaine ne signifie pas simplement remplacer une
personne par une API. Un formateur ne se contente pas de lire un texte. Il structure un
cours, choisit des exemples, respecte un programme, s'adapte au niveau des élèves,
répond aux questions, évalue la compréhension et maintient une cohérence entre les
différents supports. Si l'on souhaite automatiser cette fonction, il faut donc automatiser
une chaîne complète de tâches, et non un seul geste isolé.

C'est cette prise de conscience qui donne au sujet sa dimension d'ingénierie. La question
n'est plus seulement : quel outil d'IA utiliser ? La question devient : comment organiser
une chaîne logicielle capable de produire, contrôler, diffuser et exploiter des contenus
pédagogiques générés par IA dans un contexte réel de formation ?

\subsection{Les essais de synthèse vocale et leurs limites}

La première brique d'automatisation envisagée après la rédaction assistée par IA a été la
synthèse vocale. L'idée était directe : si le cours peut être écrit à l'avance, alors il peut
être lu par une voix artificielle. Des outils de synthèse vocale avancés, comme
ElevenLabs, permettent aujourd'hui de produire des voix beaucoup plus naturelles que les
systèmes de lecture automatique classiques. Pour une plateforme de formation audio, ce
type de solution semble donc particulièrement attractif.

Les premiers essais ont cependant montré que le problème était plus complexe que prévu.
Produire plusieurs heures d'audio ne revient pas à générer une courte démonstration
vocale. Les textes devaient être découpés en portions compatibles avec l'outil de synthèse,
envoyés successivement, puis les fichiers audio obtenus devaient être assemblés. Des
outils de montage audio manuel, comme Audacity, pouvaient être utilisés pour regrouper
les segments, corriger les coupures ou reconstituer un cours continu.

Ce processus posait plusieurs difficultés. La première était le temps de manipulation.
Même si la voix était générée par IA, la préparation, le découpage et l'assemblage restaient
largement manuels. La deuxième était le coût. Les outils de synthèse vocale les plus
qualitatifs peuvent devenir coûteux lorsque l'on produit des volumes importants, comme
des journées complètes de formation. La troisième était la qualité pédagogique perçue :
une voix naturelle ne suffit pas si le texte lu n'est pas adapté à l'oral, s'il est répétitif ou
s'il ne respecte pas le rythme d'un cours.

Ces essais ont donc permis de clarifier un point essentiel : la synthèse vocale ne résout
pas le problème de la formation autonome à elle seule. Elle automatise la lecture, mais
pas la conception du cours. Or, pour obtenir un audio exploitable, il faut d'abord produire
un texte adapté à l'écoute. Le texte doit être clair, progressif, calibré en durée, cohérent
avec le programme et rédigé dans un style oral. Si ces conditions ne sont pas réunies, la
qualité de la voix ne compense pas les défauts du contenu.

La synthèse vocale a donc déplacé le centre du problème vers l'amont. La vraie question
n'était plus seulement de savoir comment générer une voix, mais comment générer un
script pédagogique suffisamment fiable pour être confié à une voix automatique et diffusé
à de vrais élèves.

\subsection{Les limites du cours audio préenregistré}

Même lorsque le texte et la voix sont satisfaisants, un cours audio préenregistré présente
une limite structurelle : il ne peut pas interagir avec les apprenants. Dans un cours en
direct, l'élève peut poser une question, demander une reformulation, signaler une
incompréhension ou obtenir un exemple supplémentaire. Cette interaction fait partie de
la valeur pédagogique d'un formateur.

Dans un modèle audio, cette interaction disparaît. Les élèves peuvent écouter le contenu,
mais ils ne disposent pas immédiatement d'un interlocuteur capable de répondre à leurs
questions. Cette limite est particulièrement importante dans une formation
professionnalisante, où les apprenants peuvent avoir des profils variés, des niveaux
différents et des besoins de clarification fréquents. Un passage clair pour un élève peut
être insuffisant pour un autre.

Cette difficulté a fait apparaître une deuxième dimension du problème. Il ne suffisait pas
d'automatiser la production des cours. Il fallait aussi automatiser une partie de
l'accompagnement pédagogique. Sans cela, la plateforme risquait de devenir un simple
diffuseur d'audios, capable de réduire certains coûts, mais incapable de remplacer les
fonctions d'échange et de clarification assurées par un formateur humain.

L'assistant conversationnel fondé sur la récupération d'information, ou RAG, répond à
cette contrainte. Son rôle n'est pas de discuter librement comme un chatbot généraliste,
mais de répondre aux questions des élèves à partir du contenu réel de la formation. Cette
distinction est fondamentale. Un assistant pédagogique ne doit pas inventer une réponse
plausible ; il doit s'appuyer sur les documents, les cours et les supports disponibles. Il
doit donc être contextualisé.

Cette brique d'accompagnement transforme la plateforme. Elle ne se contente plus de
diffuser un contenu produit à l'avance. Elle devient capable de proposer une réponse aux
questions des apprenants, de réduire les blocages et de prolonger l'expérience du cours.
L'enjeu n'est pas seulement fonctionnel : il touche directement à la qualité de la formation.

\subsection{D'un outil de diffusion à une plateforme autonome}

À partir de ces constats, le périmètre du projet s'est élargi. La plateforme ne devait plus
être seulement un outil de dépôt et de diffusion d'audios. Elle devait devenir une chaîne
complète de formation, capable de couvrir plusieurs étapes qui étaient auparavant
réalisées par des humains.

La première étape concerne la génération du contenu pédagogique. À partir d'un
programme, d'un référentiel ou d'un objectif de formation, la plateforme doit produire
des cours structurés. Ces cours doivent être suffisamment détaillés pour couvrir les
notions attendues, mais suffisamment maîtrisés pour ne pas dériver hors du cadre prévu.

La deuxième étape concerne la transformation de ces contenus en audios exploitables.
Cela suppose de calibrer la longueur du texte, de l'adapter à une lecture orale, de produire
les fichiers audio et de les intégrer à la journée de formation sans manipulation manuelle
excessive.

La troisième étape concerne les supports complémentaires. Un cours audio seul peut être
difficile à suivre sur la durée. Des supports visuels, comme des slides, permettent
d'accompagner l'écoute, de structurer les idées principales et de fournir des repères aux
apprenants. De la même manière, des exercices ou des QCM permettent de vérifier la
compréhension et de produire des indicateurs de progression.

La quatrième étape concerne l'accompagnement. Un assistant pédagogique doit permettre
aux élèves de poser des questions et d'obtenir des réponses cohérentes avec les contenus
enseignés. Cette brique est nécessaire pour compenser l'absence d'un formateur humain
en direct.

La cinquième étape concerne le suivi. Dans une formation réelle, la plateforme doit
permettre de suivre la présence, la progression et les interactions des élèves. Elle doit
donc produire des traces exploitables, à la fois pour l'organisation interne et pour
l'évaluation du dispositif.

Le terme de plateforme autonome désigne donc ici une plateforme capable d'automatiser
une part significative de la production, de la diffusion et de l'accompagnement pédagogique.
Il ne signifie pas que toute responsabilité humaine disparaît. Une supervision reste
nécessaire, notamment pour valider les contenus, corriger les erreurs, interpréter les
résultats et décider des évolutions pédagogiques. Mais l'objectif est bien de réduire la
dépendance aux interventions humaines répétitives, afin de rendre le modèle viable pour
une petite structure.

Le tableau suivant synthétise cette évolution. Il ne présente pas encore la solution
technique détaillée, qui sera étudiée dans la partie consacrée à l'approche proposée, mais
il permet de faire apparaître la logique du problème : chaque réponse apportée à une
contrainte économique ou organisationnelle a fait émerger une nouvelle limite.

\begin{table}[h]
\centering
\small
\begin{tabular}{|p{3.1cm}|p{4cm}|p{4cm}|p{3.2cm}|}
\hline
\textbf{Modèle envisagé} & \textbf{Principe} & \textbf{Limite principale} & \textbf{Conséquence pour le projet} \\
\hline
Formateur freelance en direct &
Un formateur anime la formation, répond aux questions et adapte son discours. &
Coût récurrent trop important pour une petite structure cherchant à développer une offre externe. &
Besoin de réduire la dépendance à l'intervention humaine continue. \\
\hline
Cours audio préparé à l'avance &
Les cours sont écrits, enregistrés puis diffusés automatiquement le jour de la formation. &
Production encore manuelle : rédaction, lecture, dépôt, contrôle et mise à jour. &
Besoin d'automatiser la chaîne de production, pas seulement la diffusion. \\
\hline
Synthèse vocale isolée &
Un outil comme ElevenLabs lit les textes préparés en amont ; le montage peut être fait avec un outil comme Audacity. &
Découpage, assemblage et coût restent problématiques ; la qualité dépend surtout du texte source. &
Besoin de générer des scripts pédagogiques calibrés pour l'oral et la durée. \\
\hline
Plateforme autonome alimentée par IA &
La plateforme génère les cours, les transforme en audio, produit des supports et accompagne les élèves. &
Risque de perte de contrôle : erreurs, dérives, manque de traçabilité, qualité difficile à mesurer. &
Besoin d'une pipeline structurée, auditable et évaluée par des indicateurs. \\
\hline
\end{tabular}
\caption{Évolution du problème, de la formation humaine vers une plateforme autonome}
\end{table}

\subsection{Une classe de problèmes au-delà du cas d'Audit Énergie}

Le cas d'Audit Énergie est spécifique par son histoire, son activité et ses contraintes.
Cependant, le problème posé dépasse le seul cadre de cette entreprise. De nombreuses
petites structures rencontrent aujourd'hui une tension similaire. Elles souhaitent proposer
des services plus riches, plus personnalisés et plus accessibles, mais elles ne disposent
pas toujours des moyens humains nécessaires pour les produire manuellement.

L'intelligence artificielle générative apparaît alors comme une solution possible. Elle peut
rédiger, reformuler, résumer, synthétiser, produire de l'audio, générer des supports ou
répondre à des questions. Mais son intégration dans un contexte professionnel ne peut
pas reposer uniquement sur des usages ponctuels. Un collaborateur qui demande à un
modèle de langage de rédiger un cours obtient parfois un résultat utile, mais ce résultat
n'est pas forcément traçable, reproductible, contrôlé ni directement exploitable en
production.

La classe de problèmes étudiée peut donc être formulée ainsi : comment transformer des
capacités d'intelligence artificielle générative, souvent utilisées de manière ponctuelle,
en une chaîne logicielle fiable, auditable et intégrée à un processus métier réel ?

Dans le cas de la formation, cette question prend une forme particulière. Il ne s'agit pas
seulement de produire du texte, mais de produire un objet pédagogique complet. Un cours
doit respecter un programme, s'adapter au niveau des apprenants, être compréhensible,
être compatible avec une durée, pouvoir être transformé en audio, donner lieu à des
supports visuels, permettre une évaluation et s'inscrire dans un parcours plus large.

Cette classe de problèmes est d'autant plus importante que l'IA générative peut donner
une illusion de simplicité. Un premier résultat convaincant peut masquer les difficultés
liées à la répétabilité, au contrôle qualité, aux coûts, aux erreurs et à l'évaluation. Pour
un prototype, il suffit parfois qu'un texte semble correct. Pour une plateforme utilisée
avec de vrais élèves, cette exigence est insuffisante. Il faut pouvoir expliquer comment le
contenu a été produit, quelles règles ont été appliquées, quelles corrections ont été
réalisées et comment la qualité est mesurée.

\subsection{Les enjeux de qualité pédagogique}

La qualité pédagogique constitue l'un des principaux enjeux du projet. Dans le contexte
d'une formation professionnelle, un contenu n'est pas seulement jugé sur sa fluidité ou
son apparence. Il doit permettre à l'apprenant de progresser vers un objectif précis. Un
bon cours doit donc être clair, structuré, fidèle au programme et adapté au niveau visé.

L'utilisation de l'IA générative rend cet enjeu plus complexe. Les modèles de langage sont
capables de produire des textes très convaincants, mais ils peuvent aussi générer des
formulations imprécises, des répétitions, des exemples mal choisis ou des informations
non vérifiées. Dans une formation, ces défauts peuvent avoir des conséquences concrètes :
un élève peut retenir une explication inexacte, mal comprendre une notion ou perdre du
temps sur un contenu qui ne correspond pas au programme.

La qualité ne peut donc pas être traitée uniquement par une relecture finale. Elle doit
être intégrée dans la chaîne de production. Il faut définir ce qu'est un cours acceptable,
quelles contraintes doivent être respectées et quels types d'erreurs doivent être détectés.
Dans le projet, cela conduit à distinguer plusieurs dimensions de qualité : la cohérence
avec le plan pédagogique, l'adaptation à l'oral, le respect du budget de durée, la
conformité des formulations, la cohérence entre l'audio, les slides et les exercices, ainsi
que la capacité de l'assistant à répondre à partir des bons documents.

Cette exigence de qualité justifie le passage d'une génération simple à une pipeline plus
structurée. Un gros prompt peut produire un résultat ponctuellement satisfaisant, mais il
ne donne pas toujours les moyens de comprendre, corriger et mesurer les erreurs. Une
pipeline découpée en étapes permet au contraire d'isoler les responsabilités : planifier,
générer, relire, corriger, transformer en audio, produire des supports et conserver les
traces de chaque étape.

\subsection{Les enjeux de traçabilité et d'auditabilité}

La traçabilité est un autre enjeu central. Lorsqu'un humain produit un cours, il peut
expliquer ses choix, citer ses sources ou justifier une progression pédagogique. Lorsqu'une
IA génère un contenu, cette justification n'est pas spontanément disponible. Si la
plateforme se contente de stocker le texte final, il devient difficile de savoir comment ce
texte a été obtenu.

Cette difficulté est importante pour plusieurs raisons. D'abord, elle complique la
correction. Si un problème apparaît dans un cours, il faut savoir s'il vient du programme
de départ, du plan généré, du texte d'une section, d'une étape de réécriture, de la
calibration audio ou d'une modification finale. Sans artefacts intermédiaires, cette analyse
devient laborieuse.

Ensuite, elle limite la confiance dans le système. Une plateforme autonome doit produire
des contenus, mais elle doit aussi permettre à un responsable pédagogique ou technique
de comprendre ce qu'elle a produit. La confiance ne vient pas seulement de la qualité
perçue du résultat final ; elle vient aussi de la capacité à contrôler le processus.

Enfin, la traçabilité est nécessaire pour l'amélioration continue. Si l'on veut mesurer les
erreurs, comparer plusieurs versions de la pipeline ou identifier les étapes les plus
fragiles, il faut conserver des données sur le fonctionnement du système. L'auditabilité
devient alors une condition de l'évaluation.

Dans ce mémoire, cet enjeu justifie l'importance accordée aux artefacts intermédiaires :
plans structurés, sections générées, versions corrigées, éléments de review, supports
associés et indicateurs de génération. Ces éléments permettent de transformer une sortie
IA opaque en un processus analysable.

\subsection{Les enjeux économiques et organisationnels}

Le projet doit également être compris à travers ses enjeux économiques. L'objectif n'est
pas seulement de construire une plateforme techniquement intéressante. Il s'agit de rendre
possible un modèle de formation compatible avec les moyens d'une petite entreprise. La
réduction des interventions humaines récurrentes est donc un objectif assumé.

Cette réduction ne concerne pas uniquement le coût direct d'un formateur. Elle concerne
aussi le temps passé à préparer les contenus, à les convertir en audio, à créer des supports,
à organiser la diffusion, à répondre aux questions répétitives et à suivre les élèves. Chaque
tâche prise isolément peut sembler limitée, mais leur accumulation rend la production
d'une formation coûteuse.

L'automatisation vise donc à agir sur plusieurs leviers. Elle doit réduire le temps de
production d'un cours, limiter les manipulations manuelles, rendre les contenus plus
faciles à mettre à jour, permettre la réutilisation de certains éléments et offrir un
accompagnement de premier niveau aux apprenants. Ce sont ces gains combinés qui
peuvent rendre le modèle plus viable.

Cependant, l'automatisation introduit aussi de nouveaux coûts. Les services d'IA ont un
prix, les erreurs doivent être contrôlées, les prompts doivent être maintenus, la plateforme
doit être surveillée et les contenus doivent être validés. Il ne suffit donc pas de comparer
un formateur humain à un modèle d'IA. Il faut comparer deux organisations : un modèle
centré sur l'intervention humaine continue et un modèle centré sur une plateforme
automatisée supervisée.

Ce point est important pour la suite du mémoire. L'évaluation du projet ne devra pas se
limiter à constater que la plateforme fonctionne. Elle devra interroger le coût réel, le
temps gagné, la qualité obtenue et les limites du dispositif. C'est à cette condition que
l'automatisation pourra être analysée comme une solution d'ingénierie, et non comme une
simple expérimentation technologique.

\subsection{Les enjeux d'évaluation}

L'un des risques d'un projet fondé sur l'intelligence artificielle est de s'arrêter à une
évaluation subjective. Un cours peut sembler correct à la lecture, une voix peut sembler
naturelle, un assistant peut donner une réponse convaincante. Mais ces impressions ne
suffisent pas pour démontrer la valeur du système.

Dans le cadre de ce mémoire, l'évaluation doit donc porter sur plusieurs dimensions. La
première est l'efficacité opérationnelle : combien de temps faut-il pour produire une
journée de formation ? Combien d'étapes restent manuelles ? Combien de personnes
sont nécessaires ? Quelle part du processus est réellement automatisée ?

La deuxième dimension est la qualité pédagogique. Elle peut être observée à travers le
respect du programme, la cohérence du plan, la clarté des contenus, le taux de corrections
nécessaires, la qualité des supports visuels ou encore les résultats aux QCM. Même si
toutes ces mesures ne sont pas toujours disponibles de manière exhaustive, elles doivent
servir de guide pour analyser le système.

La troisième dimension est la qualité audio. Il faut vérifier que les textes générés peuvent
être lus naturellement, que leur durée est compatible avec les créneaux prévus et que la
voix reste acceptable pour une écoute prolongée. Le respect du budget de mots et la
correspondance entre texte et durée audio deviennent alors des indicateurs importants.

La quatrième dimension concerne l'assistant pédagogique. Un assistant RAG doit être
évalué sur la pertinence des passages retrouvés, la fidélité de ses réponses au contenu
source, sa capacité à éviter les hallucinations et son utilité réelle pour les apprenants. Une
réponse fluide mais non fondée sur les documents serait problématique.

La cinquième dimension concerne l'expérience utilisateur. La plateforme est utilisée avec
de vrais élèves ; il faut donc tenir compte de leur capacité à suivre les cours, à utiliser les
supports, à poser des questions et à progresser dans le parcours. Des retours qualitatifs,
des taux de complétion ou des résultats d'exercices peuvent contribuer à cette analyse.

Ces indicateurs ne sont pas seulement des résultats à présenter en fin de mémoire. Ils
orientent la manière de concevoir le système. Si l'on veut mesurer le respect du plan, il
faut produire un plan explicite. Si l'on veut mesurer les corrections, il faut conserver les
versions intermédiaires. Si l'on veut mesurer la pertinence du RAG, il faut pouvoir
identifier les sources utilisées. L'évaluation influence donc l'architecture.

\subsection{Formulation de la problématique}

À partir de ce contexte, le problème posé par le mémoire peut être formulé à deux
niveaux. Le premier niveau est celui d'Audit Énergie. L'entreprise souhaite développer
une activité de formation externe, mais elle doit le faire avec des moyens humains
limités. Le modèle traditionnel du formateur en direct est coûteux, tandis que les cours
audio préenregistrés réduisent l'interaction pédagogique. L'entreprise a donc besoin
d'une plateforme capable d'automatiser la production, la diffusion et une partie de
l'accompagnement des formations.

Le second niveau est plus général. Il concerne les petites structures qui cherchent à
industrialiser une activité de formation ou de service grâce à l'intelligence artificielle.
Elles peuvent accéder à des outils puissants, mais doivent encore résoudre des problèmes
de qualité, de traçabilité, de coût, d'intégration et d'évaluation. Le véritable enjeu n'est
pas seulement d'utiliser l'IA, mais de l'encadrer dans une architecture qui permette un
usage fiable en production.

La problématique du mémoire peut donc être formulée ainsi :

\begin{quote}
Comment concevoir une plateforme de formation autonome alimentée par
l'intelligence artificielle, capable de remplacer une part significative des interventions
humaines nécessaires à la production, à la diffusion et à l'accompagnement des cours,
tout en garantissant la qualité pédagogique, la traçabilité des contenus et la pertinence
des réponses apportées aux apprenants ?
\end{quote}

Cette formulation met en évidence les tensions principales du projet. D'un côté,
l'entreprise cherche à automatiser fortement son modèle de formation pour le rendre
économiquement viable. De l'autre, cette automatisation ne peut pas se faire au prix
d'une perte de qualité, d'une absence de contrôle ou d'une expérience dégradée pour les
élèves. Le problème d'ingénierie consiste donc à concilier autonomie et maîtrise.

\subsection{Objectifs de la solution attendue}

La problématique se décline en plusieurs objectifs opérationnels. Le premier objectif est
d'automatiser la génération des cours. La plateforme doit pouvoir produire des contenus
à partir d'un cadre pédagogique donné, en respectant une progression et un niveau
attendus. Cette génération ne doit pas être une simple rédaction libre : elle doit être
structurée et contrôlable.

Le deuxième objectif est d'automatiser la transformation audio. Les contenus générés
doivent être adaptés à une lecture vocale, calibrés en durée et envoyés vers un service de
synthèse vocale sans nécessiter un montage manuel lourd. L'objectif est de sortir d'un
modèle dans lequel chaque fichier audio demande une manipulation individuelle.

Le troisième objectif est de produire des supports complémentaires. Les slides et les
exercices ne doivent pas être ajoutés après coup comme des éléments indépendants. Ils
doivent être cohérents avec le contenu du cours et participer à l'expérience pédagogique.

Le quatrième objectif est d'assurer un accompagnement des apprenants. L'assistant
conversationnel doit permettre aux élèves de poser des questions et de recevoir des
réponses contextualisées. Il doit donc s'appuyer sur les contenus de formation plutôt que
sur une connaissance générale non contrôlée.

Le cinquième objectif est de rendre le système auditable. Les contenus générés doivent
être associés à des traces : plan, sections, corrections, sources, supports et indicateurs. Ce
point est nécessaire pour comprendre les erreurs, améliorer la pipeline et justifier les
choix réalisés.

Le sixième objectif est de définir des indicateurs d'évaluation. L'automatisation doit être
mesurable à travers le temps de production, le nombre d'interventions humaines, le coût,
la qualité des contenus, le respect des durées, la satisfaction des élèves, les résultats aux
exercices et la pertinence des réponses de l'assistant.

\subsection{Motivation personnelle et intérêt académique du sujet}

Ce sujet présente un intérêt particulier dans le cadre d'un mémoire d'ingénieur, car il se
situe à la frontière entre un besoin métier concret et des problématiques techniques
actuelles. Il ne s'agit pas d'étudier l'intelligence artificielle de manière abstraite, mais de
l'intégrer dans un système utilisé par une entreprise, avec de vrais élèves, des contraintes
de coût, des attentes de qualité et des limites opérationnelles.

La motivation du travail vient également de la progression du projet. La plateforme n'est
pas née directement sous sa forme actuelle. Elle a évolué par étapes : diffusion d'audios,
préparation manuelle des cours, essais de synthèse vocale, réflexion sur l'automatisation
de bout en bout, puis structuration d'une pipeline de génération et ajout d'un assistant
RAG. Cette évolution rend le sujet intéressant, car chaque étape a fait apparaître une
nouvelle limite et a obligé à reformuler le problème.

Ce mémoire permet donc de prendre du recul sur une réalisation technique. La question
n'est pas seulement de montrer ce qui a été développé, mais d'expliquer pourquoi ces
choix ont été faits, quelles alternatives auraient été possibles, quels compromis ont été
acceptés et comment les résultats peuvent être évalués. Cette prise de recul est
nécessaire pour éviter de présenter la plateforme comme une simple accumulation
d'outils d'IA.

L'intérêt académique du sujet tient aussi à la tension entre automatisation et qualité. Les
modèles génératifs rendent possible une production rapide de contenus, mais leur usage
professionnel exige des mécanismes de contrôle. Dans le domaine de la formation, cette
exigence est encore plus forte, car le contenu produit sert directement à l'apprentissage
d'élèves. La solution doit donc être pensée comme une architecture contrôlée, et non
comme une succession de générations isolées.

\subsection{Délimitation du périmètre}

Le périmètre du mémoire ne couvre pas l'ensemble des problématiques possibles de la
formation professionnelle. Il se concentre sur la conception d'une plateforme autonome
pour la production, la diffusion et l'accompagnement de cours à distance. Les aspects
administratifs, commerciaux ou réglementaires de l'organisme de formation ne sont
abordés que lorsqu'ils éclairent les contraintes du projet.

De même, le mémoire ne vise pas à démontrer que l'IA peut remplacer toute forme
d'enseignement humain dans tous les contextes. Le cas étudié est plus précis : il s'agit de
formations à distance, structurées, répétables, pour lesquelles une partie importante de la
production et de l'accompagnement peut être automatisée. La supervision humaine reste
nécessaire, mais elle change de rôle. Elle intervient davantage dans la validation, le
pilotage, l'amélioration et le traitement des cas qui dépassent le périmètre prévu.

Enfin, le mémoire ne se limite pas à la performance technique des modèles utilisés. Les
modèles de langage, les outils de synthèse vocale et les systèmes RAG sont des briques
importantes, mais le problème principal est leur orchestration dans une plateforme fiable.
L'étude porte donc sur l'architecture, les choix de pipeline, les mécanismes de contrôle et
les indicateurs permettant d'évaluer l'apport de l'automatisation.

\subsection{Transition vers l'état de l'art}

La position du problème fait apparaître plusieurs familles de solutions à comparer. Les
plateformes de formation en ligne classiques permettent de diffuser des contenus et de
suivre des apprenants, mais elles ne produisent généralement pas automatiquement des
cours complets. Les outils de synthèse vocale automatisent la lecture, mais pas la
conception pédagogique. Les modèles de langage peuvent générer du texte, mais ils
doivent être encadrés pour éviter les dérives, les répétitions et les erreurs. Les systèmes
RAG permettent de contextualiser les réponses, mais leur qualité dépend de la recherche
documentaire et de la fidélité aux sources.

L'état de l'art devra donc analyser ces approches séparément avant de les replacer dans
une même problématique : comment passer d'outils puissants mais isolés à une
plateforme de formation autonome, intégrée et contrôlable ? Cette comparaison permettra
ensuite de situer la solution développée dans le mémoire, d'en montrer l'originalité et
d'en discuter les limites.
